\documentclass[11pt,aspectratio=169]{ctexbeamer}

\usetheme{default}
\setbeamertemplate{navigation symbols}{}

\usepackage{amsmath}
\usepackage{booktabs}

% Projector-safe palette: high contrast on common lecture-room projectors.
\definecolor{DeepBlue}{RGB}{11,58,102}
\definecolor{SoftInk}{RGB}{31,41,51}
\definecolor{PaperBg}{RGB}{255,255,255}
\definecolor{PanelBlue}{RGB}{226,237,247}

\setbeamercolor{normal text}{fg=SoftInk,bg=PaperBg}
\setbeamercolor{structure}{fg=DeepBlue}
\setbeamercolor{frametitle}{fg=DeepBlue,bg=PaperBg}
\setbeamercolor{title}{fg=DeepBlue}
\setbeamercolor{block title}{fg=DeepBlue,bg=PanelBlue}
\setbeamercolor{block body}{fg=SoftInk,bg=white}

\newcommand{\Core}[1]{
  \begin{block}{核心结论}
    #1
  \end{block}
}

\newcommand{\LectureMeta}[2]{
  \title{大模型推理系统与实践}
  \subtitle{#1}
  \author{张书豪}
  \institute{华中科技大学计算机学院}
  \date{2027 · #2}
}

\newcommand{\PaperRef}[3]{
  \item \textbf{#1}（#2）\hfill\href{#3}{原文链接}
}

\newcommand{\ResearchSources}[1]{
  \begin{frame}{本讲参考资料}
    \normalsize
    \begin{itemize}
      #1
    \end{itemize}
    \vspace{0.3em}
    \textcolor{SoftInk}{课件中的论文数字均指作者在特定模型、硬件和工作负载下报告的结果，不代表无条件性能。}
  \end{frame}
}

\newcommand{\InternalCaseNote}{
  \small\textcolor{SoftInk}{说明：本页为课程研究案例或教学假设，不代表已发表论文结论；只有在实验包给出模型、硬件、工作负载和原始记录时，数字才可作为结果引用。}
}

% 14 讲细分版的统一参考资料入口。论文数字的解释边界见课程内容基线。

\newcommand{\LectureOneSources}{\ResearchSources{
  \PaperRef{vLLM: Easy, Fast, and Cheap LLM Serving with PagedAttention}{SOSP 2023}{https://doi.org/10.1145/3600006.3613165}
  \PaperRef{SGLang: Efficient Execution of Structured Language Model Programs}{arXiv:2312.07104}{https://arxiv.org/abs/2312.07104}
  \PaperRef{nano-vllm-hust 教学代码}{课程仓库}{https://github.com/intellistream/intro-to-llm-inference-engines}
}}

\newcommand{\LectureTwoSources}{\ResearchSources{
  \PaperRef{vLLM / PagedAttention}{SOSP 2023}{https://doi.org/10.1145/3600006.3613165}
  \PaperRef{Sarathi-Serve}{OSDI 2024}{https://www.usenix.org/conference/osdi24/presentation/agrawal}
  \PaperRef{DistServe}{OSDI 2024}{https://www.usenix.org/conference/osdi24/presentation/zhong-yinmin}
}}

\newcommand{\LectureThreeSources}{\ResearchSources{
  \PaperRef{vLLM / PagedAttention}{SOSP 2023}{https://doi.org/10.1145/3600006.3613165}
  \PaperRef{DistServe}{OSDI 2024}{https://www.usenix.org/conference/osdi24/presentation/zhong-yinmin}
  \PaperRef{Llumnix}{OSDI 2024}{https://www.usenix.org/conference/osdi24/presentation/sun-biao}
}}

\newcommand{\LectureFourSources}{\ResearchSources{
  \PaperRef{Orca}{OSDI 2022}{https://www.usenix.org/conference/osdi22/presentation/yu}
  \PaperRef{Sarathi-Serve}{OSDI 2024}{https://www.usenix.org/conference/osdi24/presentation/agrawal}
  \PaperRef{FastServe}{NSDI 2026}{https://www.usenix.org/conference/nsdi26/presentation/wu-bingyang}
}}

\newcommand{\LectureFiveSources}{\ResearchSources{
  \PaperRef{vLLM / PagedAttention}{SOSP 2023}{https://doi.org/10.1145/3600006.3613165}
  \PaperRef{vAttention}{arXiv:2405.04437}{https://arxiv.org/abs/2405.04437}
  \PaperRef{Rethinking KV Cache Compression}{MLSys 2025}{https://proceedings.mlsys.org/paper_files/paper/2025/hash/26289c647c6828e862e271ca3c490486-Abstract-Conference.html}
  \PaperRef{Mooncake}{FAST 2025}{https://www.usenix.org/conference/fast25/presentation/qin}
}}

\newcommand{\LectureSixSources}{\ResearchSources{
  \PaperRef{DistServe}{OSDI 2024}{https://www.usenix.org/conference/osdi24/presentation/zhong-yinmin}
  \PaperRef{RAGCache}{arXiv:2404.12457}{https://arxiv.org/abs/2404.12457}
  \PaperRef{Mooncake}{FAST 2025}{https://www.usenix.org/conference/fast25/presentation/qin}
  \PaperRef{NEO}{MLSys 2025}{https://proceedings.mlsys.org/paper_files/paper/2025/hash/66a026c0d17040889b50f0dfa650e5e0-Abstract-Conference.html}
}}

\newcommand{\LectureSevenSources}{\ResearchSources{
  \PaperRef{vLLM / PagedAttention}{SOSP 2023}{https://doi.org/10.1145/3600006.3613165}
  \PaperRef{SGLang}{arXiv:2312.07104}{https://arxiv.org/abs/2312.07104}
  \PaperRef{DistServe}{OSDI 2024}{https://www.usenix.org/conference/osdi24/presentation/zhong-yinmin}
  \PaperRef{Mooncake}{FAST 2025}{https://www.usenix.org/conference/fast25/presentation/qin}
}}

\newcommand{\LectureEightSources}{\ResearchSources{
  \PaperRef{FlashAttention}{NeurIPS 2022}{https://arxiv.org/abs/2205.14135}
  \PaperRef{FlexGen}{ICML 2023}{https://proceedings.mlr.press/v202/sheng23a.html}
  \PaperRef{PowerInfer}{SOSP 2024 / arXiv:2312.12456}{https://arxiv.org/abs/2312.12456}
  \PaperRef{FlashInfer}{MLSys 2025}{https://proceedings.mlsys.org/paper_files/paper/2025/hash/dbf02b21d77409a2db30e56866a8ab3a-Abstract-Conference.html}
}}

\newcommand{\LectureNineSources}{\ResearchSources{
  \PaperRef{DEEPSERVE}{USENIX ATC 2025}{https://www.usenix.org/conference/atc25/presentation/hu-junhao}
  \PaperRef{QServe}{MLSys 2025}{https://proceedings.mlsys.org/paper_files/paper/2025/hash/fbe2b2f74a2ece8070d8fb073717bda6-Abstract-Conference.html}
  \PaperRef{vLLM-HUST}{开源维护仓库}{https://github.com/vLLM-HUST/vllm-hust}
}}

\newcommand{\LectureTenSources}{\ResearchSources{
  \PaperRef{vLLM / PagedAttention}{SOSP 2023}{https://doi.org/10.1145/3600006.3613165}
  \PaperRef{Revisiting Disaggregated LLM Serving}{2025 preprint, arXiv:2601.08833}{https://arxiv.org/abs/2601.08833}
  \PaperRef{Rethinking KV Cache Compression}{MLSys 2025}{https://proceedings.mlsys.org/paper_files/paper/2025/hash/26289c647c6828e862e271ca3c490486-Abstract-Conference.html}
}}

\newcommand{\LectureElevenSources}{\ResearchSources{
  \PaperRef{MLPerf Inference}{MLCommons benchmark}{https://mlcommons.org/benchmarks/inference-datacenter/}
  \PaperRef{ACM Artifact Review and Badging}{ACM policy}{https://www.acm.org/publications/policies/artifact-review-and-badging-current}
  \PaperRef{Revisiting Disaggregated LLM Serving}{2025 preprint, arXiv:2601.08833}{https://arxiv.org/abs/2601.08833}
}}

\newcommand{\LectureTwelveSources}{\ResearchSources{
  \PaperRef{nano-vllm-hust}{课程教学代码}{https://github.com/intellistream/intro-to-llm-inference-engines}
  \PaperRef{vLLM-HUST}{HUST 维护的 vLLM fork}{https://github.com/vLLM-HUST/vllm-hust}
  \PaperRef{上游 vLLM}{开源项目}{https://github.com/vllm-project/vllm}
}}

\newcommand{\LectureThirteenSources}{\ResearchSources{
  \PaperRef{课程内容基线与实验材料}{课程仓库}{https://github.com/intellistream/intro-to-llm-inference-engines}
  \PaperRef{DEEPSERVE}{USENIX ATC 2025}{https://www.usenix.org/conference/atc25/presentation/hu-junhao}
  \PaperRef{GhostServe}{MLSys 2026}{https://proceedings.mlsys.org/paper_files/paper/2026/hash/d99e8e80a6c41e148db686918dd7eab3-Abstract-Conference.html}
}}

\newcommand{\LectureFourteenSources}{\ResearchSources{
  \PaperRef{课程完整材料与后续更新}{课程仓库}{https://github.com/intellistream/intro-to-llm-inference-engines}
  \PaperRef{SpecGen}{MLSys 2026}{https://proceedings.mlsys.org/paper_files/paper/2026/hash/66a026c0d17040889b50f0dfa650e5e0-Abstract-Conference.html}
  \PaperRef{GhostServe}{MLSys 2026}{https://proceedings.mlsys.org/paper_files/paper/2026/hash/d99e8e80a6c41e148db686918dd7eab3-Abstract-Conference.html}
}}

